\FifthArticle{20}
Zur weiteren Erläuterung diene als Beispiel die Berechnung der imaginären Wurzeln der oben behandelten Gleichung
\[
x^7+28x^4-480=0
\]
Nach der Bezeichnung von Artikel 17 haben wir hier zunächst wie oben $m=4,\ n=3,\ e=28,\ f=480$ und außerdem $\varepsilon=0,\ \varphi=180^{0}$. Die Formeln I von Artikel 18 werden daher
\[
\begin{aligned}
r^4&=\frac{480\sin7\rho}{28\sin3\rho}\\
r^7&=-\frac{480\sin4\rho}{\sin3\rho}\\
r^3&=-\frac{28\sin4\rho}{\sin7\rho}
\end{aligned}
\]
und die Formel II wird
\[
\frac{1}{\lambda}=\frac{823543}{6750}=\frac{\sin7\rho^7}{\sin3\rho^3\sin4\rho^4}
\]
Aus dieser Gleichung sind zwei zwischen $0$ und $90^{0}$ liegende Werte von $\rho$ zu bestimmen, da die Gleichung $\mathrm{X}=0$ neben ihren drei bereits ermittelten reellen Wurzeln noch zwei Paare zusammengehöriger imaginärer Wurzeln hat. Innerhalb dieser Grenzen wird $\sin7\rho$ dreimal $=0$, nämlich für $\rho=25\frac57$ Grad, $51\frac37$ Grad und $77\frac17$ Grad, wobei $\sin7\rho^7$ jedes Mal sein Vorzeichen ändert. $\sin3\rho$ wird einmal $=0$, für $\rho=60^{0}$, ebenfalls mit Vorzeichenwechsel von $\sin3\rho^3$. Schließlich wird $\sin4\rho$ einmal $=0$, für $\rho=45^{0}$, aber ohne Vorzeichenwechsel von $\sin4\rho^4$. Berücksichtigt man nun noch, dass der Wert von $\frac{\sin7\rho^7}{\sin3\rho^3\sin4\rho^4}$ für $\rho=0$ gleich dem Grenzwert $\frac{7^7}{3^3 4^4}$ zu setzen ist, so lässt sich das Verhalten der Werte jener Funktion in den sechs Unterabschnitten des Intervalls von $0$ bis $90^{0}$ in folgender Übersicht zusammenfassen:
\SourceLeaf{107}
\[
\begin{array}{c|c}
\rho=0\ \text{Grad}&\frac{823543}{6912}\\
&+\\
25\frac57&0\\
&-\\
45&\infty\\
&-\\
51\frac37&0\\
&+\\
60&\infty\,.\\
&-\\
77\frac17&0\\
&+\\
90^{0}&\infty
\end{array}
\]
Daraus erkennt man, dass sowohl im vierten als auch im sechsten Intervall ein Wert von $\rho$ liegen muss, der der Formel II genügt. Mehr braucht es für unseren Zweck nicht, da schon von vornherein feststeht, dass es nur zwei solche Werte gibt. Ich bringe die Gleichung II in die Form
\[
7\log\sin7\rho-3\log\sin3\rho-4\log\sin4\rho=S=2{,}0863825
\]
Die folgende Übersicht zeigt, wie sich der zwischen $51\frac37$ und $60$ Grad liegende Wert durch schrittweise Näherung mit Tafeln von $3,\ 4,\ 5,\ 7$ Ziffern finden lässt:
\[
\begin{array}{c|c|c}
\rho&S&\text{Fehler}\\
57^{0}&1{,}527&-0{,}559\\
58&2{,}354&+0{,}268
\end{array}
\]
\[
\begin{array}{c|c|c}
57^{0}40'&2{,}0624&-0{,}0240\\
57\ 50&2{,}2057&+0{,}1193
\end{array}
\]
\[
\begin{array}{c|c|c}
57^{0}41'&2{,}07658&-0{,}00980\\
57\ 42&2{,}09074&+0{,}00436
\end{array}
\]
\[
\begin{array}{c|c|c}
57^{0}41'41''&2{,}0862962&-0{,}0000863\\
57\ 41\ 42&2{,}0865320&+0{,}0001495
\end{array}
\]
Hieraus folgt $\rho=57^{0}41'41''366$ und weiter nach der zweiten der Formeln I
\SourceLeaf{108}
\[
\begin{array}{lrcl}
\log\sin4\rho&&=&9{,}8891425\,n\\
\mathrm{Compl.}\,\log\sin3\rho&&=&0{,}9193523\\
\log(-480)&&=&2{,}6812412\,n\\\hline
7\log r&&=&3{,}4897360\\
\log r&&=&0{,}4985337
\end{array}
\]
und damit
\[
x=+1{,}6843159+2{,}6637914\,i
\]
sowie die andere zugehörige Wurzel
\[
x=+1{,}6843159-2{,}6637914\,i
\]

Der andere zwischen $77\frac{1}{7}$ und $90^{0}$ liegende Wert von $\rho$ wird mit Tafeln mit drei Dezimalstellen als zwischen $86^{0}$ und $87^{0}$ liegend erkannt. Die Rechnung hat dieselbe Form wie im vorhergehenden Fall:
\[
\begin{array}{c|r|r}
\rho&S&\mathrm{Fehler}\\
86^{0}&1{,}885&-0{,}201\\
87&2{,}533&+0{,}447\\
\end{array}
\]
\[
\begin{array}{c|r|r}
86^{0}10'&1{,}9907&-0{,}0957\\
86\;20&2{,}0946&+0{,}0082\\
\end{array}
\]
\[
\begin{array}{c|r|r}
86\;19&2{,}08409&-0{,}00229\\
86\;20&2{,}09447&+0{,}00809\\
\end{array}
\]
\[
\begin{array}{c|r|r}
86^{0}19'13''&2{,}0863229&-0{,}0000596\\
86\;19\;14&2{,}0864970&+0{,}0001145
\end{array}
\]
\[
\rho=86^{0}19'13''342
\]
\[
\begin{array}{lrcl}
\log\sin4\rho&&=&9{,}4049540\,n\\
\mathrm{Compl.}\,\log\sin3\rho&&=&0{,}0081108\,n\\
\log(-480)&&=&2{,}6812412\,n\\\hline
7\log r&&=&2{,}0943060\,n\\
\log r&&=&0{,}2991866\,n
\end{array}
\]
Zieht man es vor, $r$ positiv zu haben, so braucht man nur zugleich für $\rho$ den um $180^{0}$ vergrößerten Wert $266^{0}19'13''342$ anzusetzen. Die Wurzel selbst ist
\[
x=-0{,}1278113-1{,}9874234\,i
\]
und die andere zugehörige unterscheidet sich nur im Vorzeichen des imaginären Teils davon.
\par\SourceLeaf{109}
Sämtliche Wurzeln der Gleichung $x^7+28x^4-480=0$ sind daher
\[
\begin{array}{l}
+1{,}9228841\\
-2{,}4580892\\
-2{,}5778036\\
+1{,}6843159+2{,}6637914\,i\\
+1{,}6843159-2{,}6637914\,i\\
-0{,}1278113+1{,}9874234\,i\\
-0{,}1278113-1{,}9874234\,i
\end{array}
\]
Die Summe der Wurzeln, $+0{,}0000005$, stimmt so genau mit dem wahren Wert $0$ überein, wie es vom Gebrauch siebenstelliger Logarithmen nur erwartet werden durfte. In der anderen Form hat man
\[
\begin{array}{r|r}
\log r&\rho\\\hline
0{,}2839531&0\\
0{,}3905976&180^{0}\\
0{,}4112498&180\\
0{,}4985337&57\;41'41''366\\
0{,}4985337&302\;18\;18{,}634\\
0{,}2991866&93\;40\;46{,}658\\
0{,}2991866&266\;19\;13{,}342
\end{array}
\]
Die Summe der Logarithmen der Werte von $r$ ergibt sich zu $=2{,}6812411$ und stimmt ebenfalls mit befriedigender Genauigkeit mit dem Logarithmus von $480$ überein.

Es ist wohl kaum nötig, daran zu erinnern, dass die in diesem Artikel sowie in Artikel 16 angegebenen Rechnungen nur den Gang der Arbeit in seinen Hauptschritten erläutern sollen. Für die Form des kleinen Mechanismus der Operationen sollen sie keineswegs maßgebend sein. Wer im Rechnen geübter ist, wird meistens weniger Zwischenstufen bevorzugen als die in jenen Beispielen verwendeten. Überhaupt wird jeder, der mit solchen Arbeiten einigermaßen vertraut ist, die Einzelheiten leicht selbst in die Form bringen, die den jeweiligen Umständen und seiner eigenen Gewohnheit am besten entspricht. Hier kann nicht der Ort sein, auf solche Einzelheiten weiter einzugehen.
\begin{center}\rule{0.25\linewidth}{0.3pt}\end{center}
